\begin{table}[htbp]
  \centering
  \caption{Cost ratio, timestep-ceiling ratio and the resulting advantage for the two
  main-text configurations. The cost ratio is the measured wall time per step of the
  particle engine over that of the lattice at the same box and the same particle
  content. The ceiling ratio is the largest step each method may take at a declared
  accuracy budget, with the particle engine quoted at its Euler stability limit
  rather than at the smaller step it is configured with. Matched physics means the
  detailed-balance reaction handler, which the particle engine implements for its
  single-core kernel only; best kernel means its multithreaded kernel with the
  standard handler, which is faster and is a different reaction model. The ternary
  mixture carries a sphere larger than a subvolume, which the single-core lattice
  reference does not implement, so that system is reported on the device arm only.}
  \label{tab:advantage}
  \begin{tabular}{llrrrr}
  \toprule
  system & pairing & $c_{\mathrm{R}}/c_{\mathrm{L}}$ & $\tau_{\mathrm{L}}$ ($\mu$s) & $\tau_{\mathrm{L}}/\tau_{\mathrm{R}}$ & $A$ \\
  \midrule
  ternary mixture & GPU against one core & 899 & 53.2 & 133 & $1.2\times10^{5}$ \\
   & GPU against one socket & 127 &  &  & $1.7\times10^{4}$ \\
  \addlinespace
  sinusoidal, reversible binding & one core, matched physics & 7.37 & 53.1 & 133 & 978 \\
   & one core, best kernel & 7.66 &  &  & $1.0\times10^{3}$ \\
   & one socket each & 97.13 &  &  & $1.3\times10^{4}$ \\
   & GPU against one core & 180 &  &  & $2.4\times10^{4}$ \\
   & GPU against one socket & 146 &  &  & $1.9\times10^{4}$ \\
  \bottomrule
  \end{tabular}
\end{table}
